\documentclass[10pt,a4paper]{amsart}
\usepackage[margin=19mm]{geometry}
\usepackage{amsmath,amssymb,mathtools}
\usepackage[scr=rsfs,cal=euler]{mathalfa}
\usepackage{xfrac}
\usepackage[unicode,hidelinks,bookmarks=false]{hyperref}
\newcommand{\ie}{i.e.\ }
\newcommand{\cf}{cf.\ }
\newcommand{\resp}{respectively\ }
\newcommand{\Subsection}{\subsection}
\providecommand{\nobreakdash}{\nobreak\hbox{-}}
\setlength{\emergencystretch}{2em}
\multlinegap0pt
\usepackage{amssymb}
\usepackage{tikz}
\usepackage{float}
\newtheorem{lemma}{Lemma}
\usepackage{cleveref}
\crefname{lemma}{Lemma}{Lemmas}
\providecommand{\Bignorm}[1]{\Bigl\lVert#1\Bigr\rVert}
\providecommand{\norm}[1]{\lVert#1\rVert}
\title{A structural bound for cluster robustness of randomized small-block Lanczos}
\author{Nian Shao}
\date{Source: arXiv:2507.10144v2 (CC BY 4.0)}
\begin{document}
\maketitle

\noindent\emph{Source and licence.} A structural bound for cluster robustness of randomized small-block Lanczos. Version arXiv:2507.10144v2 (CC BY 4.0), distributed by arXiv under the Creative Commons Attribution 4.0 International licence, http://creativecommons.org/licenses/by/4.0/, as recorded in the arXiv licence metadata for this exact version and shown on the abstract page https://arxiv.org/abs/2507.10144v2. Full licence notice: \texttt{licenses/arxiv-2507.10144-CC-BY-4.0.txt}.\par\smallskip

\noindent\textit{Editorial scope.} Counted unit: the article's lemma lem:sb (source0.tex lines 1017--1027). The article's complete original proof of it is delivered with it (source0.tex lines 1029--1080, one \texttt{proof} environment of 52 lines). No mathematics is written, repaired or re-derived: the statement and the proof are the article's own lines, in the article's own notation, and no retained line is substituted or re-flowed.  No further source-local context is needed: every cross-reference of every retained line resolves inside the two delivered spans.  Nothing was omitted from any retained span, so the package carries no omission record.  No retained line contains a citation, so the wrapper carries no bibliography and no bibliography entry.  No retained line contains a figure environment, an image-inclusion command or drawing code, so no image is delivered and the package declares no asset dependency.  Editorial formatting only, in addition: an AMS \texttt{amsart} preamble that restates the article's own declarations and macros used by the retained lines, namely the environments \texttt{lemma} on separate counters, together with the standard packages those lines need.  The retained environments print in this document as Lemma 1 in order of appearance, whereas the article prints the same items with the numbering of its own full document.  The complete native source of the article as distributed in the arXiv e-print for this exact version is bundled unchanged as \texttt{sources/181533/source0.tex}.
\begin{lemma}
    \label{lem:sb}
    Suppose that $B_{1}-B_{2}$ is nonsingular, then 
    \begin{equation*}
        \frac{3-\sqrt{5}}{2}\norm{(B_{1}-B_{2})^{-1}}^{2}\leq 
        \Bignorm{\begin{bmatrix}
        I & B_{1} \\ 
        I & B_{2}
    \end{bmatrix}^{-1}}^{2} \leq \frac{3+\sqrt{5}}{2}\bigl(2+(2\norm{B_{1}}^{2}+1)\norm{(B_{1}-B_{2})^{-1}}^{2}\bigr).
    \end{equation*}
\end{lemma}

\begin{proof}
Let $X=-B_{1}$ and $Y=(B_{2}-B_{1})^{-1}$. The block LU factorization yields  
\begin{equation}
    \label{eq:appLU}
    \begin{bmatrix}
        I & B_{1} \\ 
        I & B_{2}
    \end{bmatrix}^{-1}
    = 
    \begin{bmatrix}
        I & XY \\  & Y
    \end{bmatrix}
    \begin{bmatrix}
        I & \\ -I & I
    \end{bmatrix}.
\end{equation}
We just need to bound the block upper triangular matrix in \cref{eq:appLU}.
On the one hand,
\begin{equation*}
    \norm{Y} = \max_{\norm{x}=1}\norm{Yx} \leq \max_{\norm{x}=1}\Bignorm{\begin{bmatrix}
        XYx\\ Yx
    \end{bmatrix}}  = \max_{\norm{x}=1}\Bignorm{\begin{bmatrix}
        I & XY \\ 
        & Y
    \end{bmatrix} \begin{bmatrix}
        0\\ x
    \end{bmatrix}} \leq \Bignorm{\begin{bmatrix}
        I & XY \\ 
        & Y
    \end{bmatrix} }.
\end{equation*} 
On the other hand, 
\begin{equation*}
    \begin{aligned}
        \Bignorm{\begin{bmatrix}
        I & XY \\ 
        & Y
    \end{bmatrix}}^{2} 
    &= \max_{\norm{x}^{2}+\norm{y}^{2}=1}
    \Bignorm{\begin{bmatrix}
        I & XY \\ 
        & Y
    \end{bmatrix}\begin{bmatrix}
        x\\ y
    \end{bmatrix}}^{2}
    = \max_{\norm{x}^{2}+\norm{y}^{2}=1}\norm{x+XYy}^{2}+\norm{Yy}^{2} \\ 
    & \leq \max_{\norm{x}^{2}+\norm{y}^{2}=1}  \norm{x}^{2}+ 2\norm{X}\norm{Y}\norm{x}\norm{y}+(\norm{X}^{2}+1)\norm{Y}^{2}\norm{y}^{2}
    \leq 2+(2\norm{X}^{2}+1)\norm{Y}^{2}.       
    \end{aligned}
\end{equation*}
Combining these two inequalities above with the singular values of the lower triangular matrix in \cref{eq:appLU}, we finish the proof. 
\end{proof}

\end{document}
